% ===== 原文 pp.28--39 =====
% 本块：§3（原文 pp.28--33）、§4（原文 pp.34--39）。定理类编号照原文，每节从 1 起。

% [原文 p.28]
\section{第二讲（同伦、Bott 周期性与上同调性质）}

我们假设底空间紧的复向量丛。若底空间是不交并
\[ X = \coprod X_a, \]
则
\[ K(X) = \prod K(X_a). \]
此外我们已经看到，若 $G_k = \bigcup_n G_k(\C^n)$，则
\[ \Vect_k(X) = [X, G_k] \]
（确切地说，我们指的是上述两个集合的基数相等）。

我们将研究群 $K(X)$ 的同伦解释。

设底空间 $X$ 有基点 $x_0$。我们用 $\wt K$ 表示约化 K 理论（见下面的定义 2）
\[ \wt K X = \Ker(K(X)\to K(x_0)). \]
然而 $K(*) = \Z$，同时向量丛的维数给出一个映射 $X\to\Z$，其中 $[\xi]-[\eta]\mapsto\dim(\xi)-\dim(\eta)$。

% [Tutor Note: p.28 原文把维数写成一个映射 $X\to\Z$；按紧接着的 $[\xi]-[\eta]\mapsto\dim(\xi)-\dim(\eta)$ 判断，真正定义在 K 群上的是同态 $K(X)\to\Z$。两者在 $X$ 连通时一致（每个丛的纤维维数在 $X$ 的每个连通分支上是常数，故确实给出 $X\to\Z$），正文照原文译。]

\begin{defn}{1}
称 $X$ 上的两个向量丛 $\xi$ 与 $\eta$ 是\emph{稳定同构}的，若存在整数 $N, M\in\Z$ 使得
\[ \xi\oplus\C^M_X = \eta\oplus\C^N_X. \]
\end{defn}

\begin{defn}{2}
$X$ 的\emph{约化 K 群} $\wt K X$ 定义为 $X$ 上向量丛的稳定同构类所成的集合。
\end{defn}

% [Tutor Note: p.28 定义 2 把 $\wt K X$ 称作「集合」，而本节开头又把它定义为群 $\Ker(K(X)\to K(x_0))$；严格说稳定同构类在 $\oplus$ 下只构成交换幺半群，再取 Grothendieck 群才得到群。核对：$X$ 连通时 $K(X)\cong\wt K X\oplus\Z$，约化正好去掉这个 $\Z$ 直和项，见 p.32 的 $(X,\emptyset)=X^+$。]

我们知道，对 Grassmann 流形有如下包含：
\[ G_k(\C^n)\into G_k(\C^{n+1}) \]

% [原文 p.29]
取归纳极限，得到
\[ G_k = \lim_{n\to\infty}G_k(\C^n) = \bigcup_n G_k(\C^n). \]
此外
\[ G_k(\C^n)\into G_{k+1}(\C^{n+1}) \]
于是
\[ G_{k+1} = \lim_{n\to\infty}G_{k+1}(\C^{n+1}) = \bigcup_n G_{k+1}(\C^{n+1}). \]
因此
\[ G_k\into G_{k+1}, \]
这使我们能够定义归纳极限
\[ G_\infty = \lim_{k\to\infty}G_k = \bigcup_k G_k. \]
直观上，$G_\infty$ 就是变量 $k$ 与 $n$ 都趋于无穷时 $G_k(\C^n)$ 的双重归纳极限。

于是有下述结果：
\begin{thm}{1}
\[ \wt K X = [X, G_\infty]. \]
\end{thm}

% [Tutor Note: p.29 定理 1 里的 $[X,G_\infty]$ 应理解为带基点映射的同伦类集合（$G_\infty=\bigcup_k G_k$ 有自然基点）；$G_\infty$ 是稳定复 Grassmann 流形，也就是 $BU$，它分类稳定复向量丛。核对：$[S^2,G_k(\C^n)]=\pi_2=\Z$（p.31）与 $\wt K(S^2)=\Z$ 相互印证。]

接下来我们研究空间 $X$ 的悬挂 $SX$ 上的向量丛。设 $X$ 是拓扑空间；记 $CX$ 为 $X$ 的锥，它定义为
\[ CX = \frac{X\times[0,1]}{X\times 0}. \]
于是 $X$ 的悬挂 $SX$ 定义为
\[ SX = C_+X\cup_X C_-X. \]
设 $f : X\to\GL_k(\C)$ 是映射。我们定义集合
\[ V(f) = \frac{C_+X\times\C^k\cup C_-X\times\C^k}{(x,v)\equiv(x,f(x)v)}. \]

% [原文 p.30]
则 $V(f)\to SX$ 是向量丛。

\begin{lem}{1}
设 $f_0, f_1 : X\to\GL_k(\C)$ 是同伦的映射。则
\[ V(f_0)\simeq V(f_1). \]
\end{lem}

此外还有下述重要结果：
\begin{thm}{2}
\[ \Vect_k(SX) = [X, \GL_k(\C)]. \]
\end{thm}

若考虑包含
\[ \GL_k(\C)\into\GL_{k+1}(\C) \]
它由
\[ A\mapsto\begin{pmatrix} A & 0 \\ 0 & 1 \end{pmatrix} \]
给出，并令
\[ \GL\C = \lim_{k\to\infty}\GL_k(\C) = \bigcup_k \GL_k(\C), \]
则定理 2 有下述推广：
\begin{thm}{3}
\[ \wt K(SX) = [X, \GL\C]. \]
\end{thm}

定理 2 与定理 3 的证明可参见 Atiyah 的笔记。

% [Tutor Note: p.30 原文对定理 2、3 只说「见 Atiyah 的笔记」，略去的论证是标准的粘合（clutching）构造：$f:X\to\GL_k(\C)$ 通过 $V(f)$ 给出 $SX$ 上的 $k$ 维丛；引理 1 说明同伦的 $f$ 给出同构的丛，故 $V$ 在 $\pi_0$ 层面诱导映射 $[X,\GL_k(\C)]\to\Vect_k(SX)$，其逆由沿 $S^1$ 方向的转移函数给出。]

\para{注：}
我们提到过
\[ \wt K X = [X, G_\infty] \]
而
\[ \wt K SX = [X, \GL\C]. \]
然而
\[ \wt K SX = [SX, G_\infty] = [X, \Omega G_\infty], \]
其中 $\Omega G_\infty$ 是 $G_\infty$ 的环路空间，即
\[ \Omega G_\infty = \mathrm{Maps}_*(S^1, G_\infty). \]

% [原文 p.31]
于是
\[ \Omega G_\infty \simeq \GL\C, \]
这意味着\emph{无穷 Grassmann 流形} $G_\infty$ \emph{是无穷群} $\GL\C$ \emph{的分类空间}：
\[ G_\infty = B\GL\C. \]

K 理论中的基本结果是\emph{Bott 周期性}，我们下面就来研究它。它最简单的形式是说，$K(X)\otimes K(S^2)$ 与 $K(X\times S^2)$ 之间存在同构。

我们记 $S^2X$ 为 $X$ 的二次悬挂，换言之 $S^2X = S(SX)$。归纳地可以定义空间 $X$ 的 $n$ 次悬挂 $S^nX$。回顾拓扑学中关于球面的悬挂的下述结果：
\[ S^{2n}(S^1) = S^{2n+1} \]
以及
\[ S^{2n-2}(S^2) = S^{2n}. \]
作为 Bott 周期性的第一个版本，我们指出：
\[ \wt K(S^2X) \simeq \wt K X \]
而
\[ [S^1, G_k(\C^n)] = \pi_1(G_k(\C^n)) = 0 \]
因为 Grassmann 流形是可缩的。

% [Tutor Note: p.31 「because Grassmannians are contractible」有误：$G_k(\C^n)$（$0<k<n$）是复维数 $k(n-k)>0$ 的紧复流形，不可缩；$\pi_1=0$ 的正确理由是它单连通。判断依据就在下一行：原文自己写了 $\pi_2(G_k(\C^n))=\Z$，若可缩则所有同伦群应为 0。正文照原文译，此处按原文。]

此外
\[ [S^2, G_k(\C^n)] = \pi_2(G_k(\C^n)) = \Z, \]
\[ \wt K S^1 = 0, \]
\[ \wt K S^2 = \Z, \]
\[ \wt K S^{2n+1} = 0 \]
以及
\[ \wt K S^{2n} = \Z. \]

现在我们来讨论 K 理论的上同调性质。

% [原文 p.32]
设 $X$ 紧，$A\subset X$ 是紧子空间。取定某个 $n\in\N$，我们有：

\begin{defn}{3}
K 理论的负幂定义如下：
\[ \wt K^{-n}X = \wt K(S^nX). \]
\end{defn}

对相对 K 理论，有如下定义：

\begin{defn}{4}
相对 K 群定义如下：
\[ \wt K^{-n}(X,A) = \wt K^{-n}(X/A) = \wt K(S^n(X/A)). \]
\end{defn}

此外有
\[ K^{-n}X = \wt K^{-n}(X,\emptyset), \]
这里我们定义
\[ (X,\emptyset) = X^+, \]
而 $X^+$ 表示在 $X$ 上添加一个不交的基点，即
\[ \wt K(X^+) = K(X). \]
利用 Bott 周期性，我们可以延拓上述定义而得到
\[ K^n(X,A) = K^{n-2}(X,A),\quad \forall n\in\Z. \]

总结如下：

\begin{enumerate}[label=\arabic*.]
\item 群 $K^n(X,A)$ 对任意 $n$ 都可以定义，而 $K^n(X,A)\simeq K^{n+2}(X,A)$。
\item 本质上，由于 Bott 周期性，只有两个 K 群，即 $K^0(X,A)$ 与 $K^1(X,A)$，其中 $K^0(X,A) = [X/A, G_\infty]$，而 $K^1(X,A) = [X/A, \GL\C]$。
\end{enumerate}

% [Tutor Note: p.32 定义 3、4 的负幂依赖悬挂与商空间；注意定义 4 对 $(X,A)$ 用 $X/A$，而 $K^{-n}X=\wt K^{-n}(X,\emptyset)$ 用 $X^+=X/\emptyset$，两者在 $A=\emptyset$ 时一致。核对：$K^{n}(X,A)\simeq K^{n+2}(X,A)$ 正是 Bott 周期性 $\wt K(S^2X)\simeq\wt K X$（p.31）的推论。]

给定映射 $f : (X,A)\to(Y,B)$，存在一个诱导映射，它是同态
\[ K^n(X,A)\xleftarrow{f^*} K^n(Y,B) \]

% [原文 p.33]
（这是 K 理论中方向相反的函子性。）同时，给定
\[ X/A\leftarrow X\leftarrow A \]
我们可以得到一对空间的 K 理论正合列
\[
\begin{tikzcd}
K^0(X,A) \arrow[r] & K^0X \arrow[r] & K^0A \arrow[d,"\delta"] \\
K^1A \arrow[u] & K^1X \arrow[l] & K^1(X,A) \arrow[l]
\end{tikzcd}
\]

% [Tutor Note: p.33 六项正合列的方向要看清：第一行向右、第二行向左，左端竖直箭头向上，右端是连接同态 $\delta$ 向下（把 $K^0A$ 送到 $K^1(X,A)$，上标降一次）。与 p.32 末尾的 $f^*$ 反方向函子性对照；原文 "funcoriality" 是 functoriality 的笔误。]

% [原文 p.34]
\section{第三讲（乘积与 Bott 周期性的再审视）}

设 $X$ 紧且连通。由 $X$ 上向量丛的同构类集合 $\Vect(X)$，我们定义了 $K^0X = [X, \Z\times G_\infty]$，而约化 0 次 K 群定义为
\[ \wt K^0X = \Ker(K^0X\to K^0(*)) = [X, G_\infty]. \]
此外我们看到
\[ \wt K^1X = \wt K^0(SX) = [X, \GL\C], \]
而对紧子空间 $A\subset X$ 及某个 $n\in\N$ 有
\[ \wt K^{-n}X = \wt K^0(S^nX), \]
\[ \wt K^{-n}(X,A) = \wt K^{-n}(X/A), \]
\[ K^{-n}X = K^{-n}(X,\emptyset), \]
其中 $X^+ = X/\emptyset$，并且
\[ K^n(X,A) = K^{n-2}(X,A) \]
用到 Bott 周期性。

% [Tutor Note: p.34 $K^0X=[X,\Z\times G_\infty]$ 中的 $\Z$ 因子记录虚拟丛的秩，故约化群 $\wt K^0X=\Ker(K^0X\to K^0(*))$ 把它去掉，只剩分类空间因子：核对 $K^0(*)=\Z$、$\wt K^0(*)=0$ 即得。]

我们列出群 $K^n(X,A)$ 的性质（它们由 Eilenberg--Steenroad 公理推出）：

\begin{enumerate}[label=\arabic*.]
\item 对应 $(X,A)\mapsto K^n(X,A)$ 是从拓扑空间范畴到阿贝尔群范畴的函子。
\item 同伦的映射在阿贝尔群之间诱导相同的同态。
\item 存在一对空间的 K 理论正合列（见上一讲末尾）。

% [原文 p.35]
\item 该正合列是自然的。
\item 切除成立：
\[ K^n(X,A)\simeq K^n(X/A,*). \]
\item 维数公理不成立
\[ K^n(*) = 0,\quad n\in\N,\ \textit{odd} \]
而
\[ K^n(*) = \Z,\quad n\in\N,\ \textit{even}. \]
\end{enumerate}

% [Tutor Note: p.34--35 这六条正是 Eilenberg--Steenrod 公理的 K 理论版本，唯第 6 条「维数公理」不成立：维数公理要求 $n>0$ 时 $K^n(*)=0$，这里却是 $K^{2m}(*)=\Z$、$K^{2m+1}(*)=0$（以 2 为周期）；公式中的 odd/even 即奇/偶。原文把 Steenrod 拼作 "Steenroad"，正文照原文。]

于是我们得到一个\emph{广义的（或非寻常的）}上同调理论。

给定
\[ A\to X\to X/A, \]
对空间的正合列曾是
\[
\begin{tikzcd}
K^0(X,A) \arrow[r] & K^0X \arrow[r] & K^0A \arrow[d,"\delta"] \\
K^1A \arrow[u] & K^1X \arrow[l] & K^1(X,A) \arrow[l]
\end{tikzcd}
\]

若 $A\subset X$，则
\[ SA\simeq\frac{X\cup CA}{X} \]
而
\[ SA\simeq\frac{X\cup CA}{X}\leftarrow X\cup CA\to X/A. \]
此外
\[ K^{-n-1}A = K^{-n}(SA)\simeq K^{-n}(X\cup CA,X)\to K^{-n}(X\cup CA)\simeq K^{-n}(X/A) \]
连同映射
\[ K^{-1}A\to K^0(X/A) \]
以及
\[ K^0A\to K^1(X/A). \]

% [Tutor Note: p.35 「$SA\simeq(X\cup CA)/X$」把 $A$ 的悬挂写成锥的商；正合列里的 $\delta$ 正是由「对 $(X\cup CA,X)$ 的正合列」与切除 $\wt K^{-n}(X/A)\simeq\wt K^{-n}(X\cup CA,X)$ 拼出来的。核对：$K^{-n-1}A=K^{-n}(SA)$ 直接把 p.32 定义 3、4 代入即可。]

% [原文 p.36]
正合性的证明是繁琐的。

下面我们尝试对 K 理论正合列中出现的连接映射 $\delta$ 给出几何解释。

对于 $K^0(X,A)$，我们从一个对 $(\xi,\phi)$ 出发，其中 $\xi$ 是 $X$ 上的向量丛，$\phi$ 是 $\xi/A$ 的一个平凡化。这些对之间有两种等价关系：同构与同伦；因此可以定义类 $[\xi,\phi]\in K^0(X,A)$。

考虑映射
\[ \delta : K^1A\to K^0(X,A) \]
连同某个
\[ f : A\to\GL_k(\C). \]
定义
\[ \frac{X\times\C^k\ \coprod\ CA\times\C^k}{\textit{glue via } f} \]
并得到
\[ X\cup CA\simeq X/A. \]

% [Tutor Note: p.36 连接同态 $\delta$ 的几何意义：$\wt K^1A=[A,\GL\C]$ 的元素由 $f:A\to\GL_k(\C)$ 代表，把两个乘积丛沿 $f$ 粘合（公式分母的 "glue via $f$" 即「沿 $f$ 粘合」）得到 $X\cup CA\simeq X/A$ 上的丛，这就是 $\delta$ 的像。原文 "conecting" 是 connecting 的笔误。]

现在我们来研究乘积。

设
\[ \Vect(X)\times\Vect(Y)\to\Vect(X\times Y) \]
其中
\[ (\xi,\eta)\mapsto\xi\otimes\eta \]
是外张量积。对纤维有
\[ F_{x,y}(\xi\otimes\eta) = F_x(\xi)\otimes F_y(\eta), \]
因为 $\xi\otimes\eta\to X\times Y$ 是新的向量丛。

于是有诱导映射
\[ K^n(X,A)\otimes K^m(Y,B)\to K^{n+m}(X\times Y, A\times Y\cup X\times B). \]

% [原文 p.37]
\begin{enumerate}[label=\arabic*.]
\item 将上述推广为一个映射
\[ K^0X\otimes K^0Y\to K^0(X\times Y) \]
其中
\[ ([\xi_1]-[\eta_1])\otimes([\xi_2]-[\eta_2]) := [\xi_1\otimes\xi_2]+[\eta_1\otimes\eta_2]-[\eta_1\otimes\xi_2]-[\xi_1\otimes\eta_2]. \]
回顾代数拓扑中的\emph{smash 积} $X\wedge Y$：设 $X$ 有基点 $x_0$、$Y$ 有基点 $y_0$，它定义为
\[ X\wedge Y = \frac{X\times Y}{x_0\times Y\cup X\times y_0}. \]
例如，空间 $X$ 的悬挂就是与圆周作 smash 积：
\[ S^1\wedge X\simeq SX \]
以及
\[ S^n\wedge X\simeq S^nX. \]
% [Tutor Note: p.37 外积公式是双线性的定义式：把虚拟丛写成形式差并要求对加法分配，展开 4 项即得 $[\xi_1\otimes\xi_2]+[\eta_1\otimes\eta_2]-[\eta_1\otimes\xi_2]-[\xi_1\otimes\eta_2]$；核对方法是验证它在 $[\xi_1]-[\eta_1]$ 与 $[\xi_2]-[\eta_2]$ 上分别线性。]
\item 我们集中讨论约化 K 理论：现在把映射推广为
\[ \wt K^0X\otimes\wt K^0Y\to\wt K^0(X\wedge Y). \]
我们把 smash 积定义为
\[ \wt K^0(X\wedge Y) = \Ker(K^0(X\times Y)\to K^0X\oplus K^0Y). \]
我们使用显然的映射
\[ K^0(X\times Y)\to K^0X \]
它来自
\[ X = X\times y_0\to X\times Y \]
以及
\[ K^0(X\times Y)\to K^0Y \]
它来自
\[ Y = Y\times x_0\to X\times Y. \]
\end{enumerate}

% [原文 p.38]
于是我们得到：
\[ \resizebox{0.98\textwidth}{!}{$\displaystyle \wt K^0(S^n(X/A))\otimes\wt K^0(S^m(Y/B))\to\wt K^0(S^n(X/A)\wedge S^m(Y/B)) = \wt K^0(S^{n+m}\wedge(X/A)\wedge(Y/B)),$} \]
其中
\[ \frac{X\times Y}{A\times Y\cup X\times B} = \frac{X}{A}\wedge\frac{Y}{B}. \]
作为例子，考虑特殊情形 $X = Y$、$A = B = \emptyset$。我们得到：
\[ K^nX\otimes K^mX\to K^{n+m}(X\times X)\xrightarrow{\Delta^*} K^{n+m}X \]
其中最后一个映射
\[ \Delta^* : K^{n+m}(X\times X)\to K^{n+m}X \]
由映射
\[ \Delta : X\to X\times X \]
诱导，这里
\[ \Delta(x) = (x,x). \]
于是我们有下述映射
\[ K^0X\otimes K^0X\to K^0X, \]
\[ K^0X\otimes K^1X\to K^1X, \]
以及
\[ K^1X\otimes K^1X\to K^0X, \]
而这本质上就是 Bott 周期性。

% [Tutor Note: p.38 $\Delta^*$ 由对角映射 $\Delta(x)=(x,x)$ 诱导，与上积复合即得环结构 $K^0X\otimes K^0X\to K^0X$；$K^1X\otimes K^1X\to K^0X$ 正是 Bott 周期性在乘积层面的表现（原文 "esentially" 是 essentially 的笔误）。]

我们还可以定义
\[ K^*X = K^0X\oplus K^1X \]
它是一个分次交换环。

现在我们有条件来看 Bott 周期性的第二个修订版本：

映射
\[ \wt K^0(S^2)\otimes\wt K^0X\to\wt K^0(S^2X) \]

% [原文 p.39]
是同构（回忆 $\wt K^0S^2 = \Z$）。

Bott 周期性也许是计算中最有用的工具。让我们举一些例子：
\[ \wt K^1(S^{2n+1}) = \Z = \wt K^0(S^{2n}) \]
而
\[ \wt K^1(S^{2n}) = 0 = \wt K^0(S^{2n-1}). \]
此外
\[ [S^{2n+1}, \GL\C] = \pi_{2n+1}(\GL\C) = \pi_{2n+1}(\GL_N(\C)),\ N > 2n+2), \]
而
\[ [S^{2n}, \GL\C] = 0. \]
若用 $U_N$ 记 $N\times N$ 复酉矩阵组成的群（它显然是 $\GL_N(\C))$ 的子群），则（当 $N > 2n+1$ 时）有
\[ \pi_{2n+1}(U_N) = \Z, \]
而
\[ \pi_{2n}(U_N) = 0. \]
对特殊正交群 $\SO_N$，我们有（当 $N > i+1$ 时）：
\[ \pi_i(\SO_N) = \Z/2,\quad i\equiv 1(8), \]
\[ \pi_i(\SO_N) = 0,\quad i\equiv 2(8), \]
\[ \pi_i(\SO_N) = \Z,\quad i\equiv 3(8), \]
\[ \pi_i(\SO_N) = 0,\quad i\equiv 4(8), \]
\[ \pi_i(\SO_N) = 0,\quad i\equiv 5(8), \]
\[ \pi_i(\SO_N) = 0,\quad i\equiv 6(8), \]
\[ \pi_i(\SO_N) = \Z,\quad i\equiv 7(8) \]
以及
\[ \pi_i(\SO_N) = \Z/2,\quad i\equiv 1(8). \]

% [Tutor Note: p.39 的 $\SO$ 同伦群表最后一行应为 $i\equiv 0(8)$：Bott 周期性给出 $\pi_i(\SO)=\Z/2$ 当 $i\equiv 0,1 \pmod 8$。表中已列 $1,2,\dots,7$ 各一行，原文把 $i\equiv 1(8)$ 写了两遍并漏掉 $i\equiv 0(8)$ 一行；正文照原样，正确形式见此处。]

% [Tutor Note: p.39 原文 "[S^{2n+1},\GL\C] = \pi_{2n+1}(\GL\C) = \pi_{2n+1}(\GL_N(\C)),\ N>2n+2)$" 末尾多一个右括号，同页 "$\GL_N(\C))$" 处也多一个右括号；又 "onviously"、"orhtogonal" 分别为 obviously、orthogonal 的笔误。正文一律照原文。]
